\documentclass{article}
\usepackage{amsmath, amssymb, ctex, graphicx}
\usepackage[margin=1in]{geometry}
\title{第6章-深度神经网络-所有习题}
\author{《深度学习基础》课程小组}
% \date{}

\begin{document}
\maketitle

\section*{Exercises}

\textbf{6.1 ($\star\star\star$)} Use the result (2.126) to derive an expression for the surface area $S_D$ and the volume $V_D$ of a hypersphere of unit radius in $D$ dimensions. To do this, consider the following result, which is obtained by transforming from Cartesian to polar coordinates:
\begin{equation}
\prod_{i=1}^{D} \int_{-\infty}^{\infty} e^{-x_i^2} \, dx_i = S_D \int_{0}^{\infty} e^{-r^2} r^{D-1} \, dr.
\tag{6.51}
\end{equation}
Using the gamma function, defined by
\begin{equation}
\Gamma(x) = \int_{0}^{\infty} t^{x-1} e^{-t} \, dt
\tag{6.52}
\end{equation}
together with (2.126), evaluate both sides of this equation, and hence show that
\begin{equation}
S_D = \frac{2\pi^{D/2}}{\Gamma(D/2)}.
\tag{6.53}
\end{equation}
Next, by integrating with respect to the radius from 0 to 1, show that the volume of the unit hypersphere in $D$ dimensions is given by
\begin{equation}
V_D = \frac{S_D}{D}.
\tag{6.54}
\end{equation}
Finally, use the results $\Gamma(1)=1$ and $\Gamma(3/2)=\sqrt{\pi}/2$ to show that (6.53) and (6.54) reduce to the usual expressions for $D=2$ and $D=3$.

\textbf{6.1 ($\star\star\star$)} 利用结果 (2.126) 推导 $D$ 维空间中单位半径超球面的表面积 $S_D$ 和体积 $V_D$ 的表达式。为此，请考虑以下通过将笛卡尔坐标变换为极坐标得到的结果：
\begin{equation}
\prod_{i=1}^{D} \int_{-\infty}^{\infty} e^{-x_i^2} \, dx_i = S_D \int_{0}^{\infty} e^{-r^2} r^{D-1} \, dr.
\tag{6.51}
\end{equation}
利用伽马函数的定义
\begin{equation}
\Gamma(x) = \int_{0}^{\infty} t^{x-1} e^{-t} \, dt
\tag{6.52}
\end{equation}
并结合 (2.126)，计算此方程的两边，从而证明
\begin{equation}
S_D = \frac{2\pi^{D/2}}{\Gamma(D/2)}.
\tag{6.53}
\end{equation}
接下来，通过对半径从 0 到 1 进行积分，证明 $D$ 维空间中单位超球体的体积由下式给出
\begin{equation}
V_D = \frac{S_D}{D}.
\tag{6.54}
\end{equation}
最后，利用结果 $\Gamma(1)=1$ 和 $\Gamma(3/2)=\sqrt{\pi}/2$ 来证明 (6.53) 和 (6.54) 在 $D=2$ 和 $D=3$ 时简化为通常的表达式。

\vspace{1em}

\textbf{6.2 ($\star\star\star$)} Consider a hypersphere of radius $a$ in $D$ dimensions together with the concentric hypercube of side $2a$, so that the hypersphere touches the hypercube at the centres of each of its sides. By using the results of Exercise 6.1, show that the ratio of the volume of the hypersphere to the volume of the cube is given by
\begin{equation}
\frac{\text{volume of hypersphere}}{\text{volume of cube}} = \frac{\pi^{D/2}}{D 2^{D-1} \Gamma(D/2)}.
\tag{6.55}
\end{equation}
Now make use of Stirling's formula in the form
\begin{equation}
\Gamma(x+1) \simeq (2\pi)^{1/2} e^{-x} x^{x+1/2},
\tag{6.56}
\end{equation}
which is valid for $x \gg 1$, to show that, as $D \to \infty$, the ratio (6.55) goes to zero. Show also that the distance from the centre of the hypercube to one of the corners divided by the perpendicular distance to one of the sides is $\sqrt{D}$, which therefore goes to $\infty$ as $D \to \infty$. From these results, we see that, in a space of high dimensionality, most of the volume of a cube is concentrated in the large number of corners, which themselves become very long 'spikes'!

\textbf{6.2 ($\star\star\star$)} 考虑一个 $D$ 维空间中半径为 $a$ 的超球体以及一个同心的、边长为 $2a$ 的超立方体，使得超球体在每个侧面的中心处与超立方体相切。通过利用练习 6.1 的结果，证明超球体体积与立方体体积之比为
\begin{equation}
\frac{\text{volume of hypersphere}}{\text{volume of cube}} = \frac{\pi^{D/2}}{D 2^{D-1} \Gamma(D/2)}.
\tag{6.55}
\end{equation}
现在使用斯特林公式的形式
\begin{equation}
\Gamma(x+1) \simeq (2\pi)^{1/2} e^{-x} x^{x+1/2},
\tag{6.56}
\end{equation}
该公式对 $x \gg 1$ 有效，以证明当 $D \to \infty$ 时，比值 (6.55) 趋于零。同时证明，从超立方体中心到一个顶点的距离除以到其中一个侧面的垂直距离是 $\sqrt{D}$，因此当 $D \to \infty$ 时它也趋于无穷大。从这些结果中我们看到，在高维空间中，立方体的大部分体积都集中在大量的顶点上，而这些顶点本身变成了非常长的“尖刺”！

\vspace{1em}

\textbf{6.3 ($\star\star\star$)} In this exercise, we explore the behaviour of the Gaussian distribution in high-dimensional spaces. Consider a Gaussian distribution in $D$ dimensions given by
\begin{equation}
p(\mathbf{x}) = \frac{1}{(2\pi\sigma^2)^{D/2}} \exp\left(-\frac{\|\mathbf{x}\|^2}{2\sigma^2}\right).
\tag{6.57}
\end{equation}
We wish to find the density as a function of the radius in polar coordinates in which the direction variables have been integrated out. To do this, show that the integral of the probability density over a thin shell of radius $r$ and thickness $\epsilon$, where $\epsilon \ll 1$, is given by $p(r)\epsilon$ where
\begin{equation}
p(r) = \frac{S_D r^{D-1}}{(2\pi\sigma^2)^{D/2}} \exp\left(-\frac{r^2}{2\sigma^2}\right)
\tag{6.58}
\end{equation}
where $S_D$ is the surface area of a unit hypersphere in $D$ dimensions. Show that the function $p(r)$ has a single stationary point located, for large $D$, at $\widehat{r} \simeq \sqrt{D}\sigma$. By considering $p(\widehat{r}+\epsilon)$ where $\epsilon \ll \widehat{r}$, show that for large $D$,
\begin{equation}
p(\widehat{r}+\epsilon) = p(\widehat{r}) \exp\left(-\frac{3\epsilon^2}{2\sigma^2}\right),
\tag{6.59}
\end{equation}
which shows that $\widehat{r}$ is a maximum of the radial probability density and also that $p(r)$ decays exponentially away from its maximum at $\widehat{r}$ with length scale $\sigma$. We have already seen that $\sigma \ll \widehat{r}$ for large $D$, and so we see that most of the probability mass is concentrated in a thin shell at large radius. Finally, show that the probability density $p(\mathbf{x})$ is larger at the origin than at the radius $\widehat{r}$ by a factor of $\exp(D/2)$. We therefore see that most of the probability mass in a high-dimensional Gaussian distribution is located at a different radius from the region of high probability density.

\textbf{6.3 ($\star\star\star$)} 在本练习中，我们探讨高斯分布在高维空间中的行为。考虑一个 $D$ 维的高斯分布，其形式为
\begin{equation}
p(\mathbf{x}) = \frac{1}{(2\pi\sigma^2)^{D/2}} \exp\left(-\frac{\|\mathbf{x}\|^2}{2\sigma^2}\right).
\tag{6.57}
\end{equation}
我们希望找到作为极坐标下半径函数的密度，其中方向变量已被积分掉。为此，请证明概率密度在一个半径为 $r$、厚度为 $\epsilon$（其中 $\epsilon \ll 1$）的薄壳上的积分由 $p(r)\epsilon$ 给出，其中
\begin{equation}
p(r) = \frac{S_D r^{D-1}}{(2\pi\sigma^2)^{D/2}} \exp\left(-\frac{r^2}{2\sigma^2}\right)
\tag{6.58}
\end{equation}
这里 $S_D$ 是 $D$ 维空间中单位超球面的表面积。证明函数 $p(r)$ 有一个唯一的驻点，对于大的 $D$，该点位于 $\widehat{r} \simeq \sqrt{D}\sigma$。通过考虑 $p(\widehat{r}+\epsilon)$（其中 $\epsilon \ll \widehat{r}$），证明对于大的 $D$，
\begin{equation}
p(\widehat{r}+\epsilon) = p(\widehat{r}) \exp\left(-\frac{3\epsilon^2}{2\sigma^2}\right),
\tag{6.59}
\end{equation}
这表明 $\widehat{r}$ 是径向概率密度的最大值，并且 $p(r)$ 从其最大值 $\widehat{r}$ 处以长度尺度 $\sigma$ 呈指数衰减。我们已经看到，对于大的 $D$，有 $\sigma \ll \widehat{r}$，因此我们看到大部分概率质量都集中在一个大半径处的薄壳中。最后，证明概率密度 $p(\mathbf{x})$ 在原点处比在半径 $\widehat{r}$ 处大 $\exp(D/2)$ 倍。因此我们看到，高维高斯分布中的大部分概率质量位于与高概率密度区域不同的半径处。

\vspace{1em}

\textbf{6.4 ($\star\star$)} Consider a two-layer network function of the form (6.11) in which the hidden-unit nonlinear activation functions $h(\cdot)$ are given by logistic sigmoid functions of the form
\begin{equation}
\sigma(a) = \{1 + \exp(-a)\}^{-1}.
\tag{6.60}
\end{equation}
Show that there exists an equivalent network, which computes exactly the same function, but with hidden-unit activation functions given by $\tanh(a)$ where the tanh function is defined by (6.14). Hint: first find the relation between $\sigma(a)$ and $\tanh(a)$, and then show that the parameters of the two networks differ by linear transformations.

\textbf{6.4 ($\star\star$)} 考虑一个形式如 (6.11) 的两层网络函数，其中隐藏单元的非线性激活函数 $h(\cdot)$ 由如下形式的逻辑 Sigmoid 函数给出
\begin{equation}
\sigma(a) = \{1 + \exp(-a)\}^{-1}.
\tag{6.60}
\end{equation}
证明存在一个等效的网络，它计算完全相同的函数，但其隐藏单元的激活函数由 $\tanh(a)$ 给出，其中 tanh 函数由 (6.14) 定义。提示：首先找到 $\sigma(a)$ 和 $\tanh(a)$ 之间的关系，然后证明两个网络的参数仅相差一个线性变换。

\vspace{1em}

\textbf{6.5 ($\star\star$)} The swish activation function (Ramachandran, Zoph, and Le, 2017) is defined by
\begin{equation}
h(x) = x \sigma(\beta x)
\tag{6.61}
\end{equation}
where $\sigma(x)$ is the logistic-sigmoid activation function defined by (6.13). When used in a neural network, $\beta$ can be treated as a learnable parameter. Either sketch or plot using software graphs of the swish activation function as well as its first derivative for $\beta = 0.1$, $\beta = 1.0$, and $\beta = 10$. Show that when $\beta \to \infty$, the swish function becomes the ReLU function.

\textbf{6.5 ($\star\star$)} Swish 激活函数 (Ramachandran, Zoph, and Le, 2017) 定义为
\begin{equation}
h(x) = x \sigma(\beta x)
\tag{6.61}
\end{equation}
其中 $\sigma(x)$ 是由 (6.13) 定义的逻辑 Sigmoid 激活函数。当在神经网络中使用时，$\beta$ 可以被视为一个可学习的参数。请绘制或使用软件绘制出 $\beta = 0.1$、$\beta = 1.0$ 和 $\beta = 10$ 时的 swish 激活函数及其一阶导数的图像。证明当 $\beta \to \infty$ 时，swish 函数变为 ReLU 函数。

\vspace{1em}

\textbf{6.6 ($\star$)} We saw in (5.72) that the derivative of the logistic-sigmoid activation function can be expressed in terms of the function value itself. Derive the corresponding result for the tanh activation function defined by (6.14).

\textbf{6.6 ($\star$)} 我们在 (5.72) 中看到，逻辑 Sigmoid 激活函数的导数可以用函数值本身来表示。请推导出由 (6.14) 定义的 tanh 激活函数的相应结果。

\vspace{1em}

\textbf{6.7 ($\star\star$)} Show that the softplus activation function $\zeta(a)$ given by (6.16) satisfies the properties:
\begin{align}
\zeta(a) - \zeta(-a) &= a \tag{6.62} \\
\ln \sigma(a) &= -\zeta(-a) \tag{6.63} \\
\frac{d\zeta(a)}{da} &= \sigma(a) \tag{6.64} \\
\zeta^{-1}(a) &= \ln(\exp(a)-1) \tag{6.65}
\end{align}
where $\sigma(a)$ is the logistic-sigmoid activation function given by (6.13).

\textbf{6.7 ($\star\star$)} 证明由 (6.16) 给出的 softplus 激活函数 $\zeta(a)$ 满足以下性质：
\begin{align}
\zeta(a) - \zeta(-a) &= a \tag{6.62} \\
\ln \sigma(a) &= -\zeta(-a) \tag{6.63} \\
\frac{d\zeta(a)}{da} &= \sigma(a) \tag{6.64} \\
\zeta^{-1}(a) &= \ln(\exp(a)-1) \tag{6.65}
\end{align}
其中 $\sigma(a)$ 是由 (6.13) 给出的逻辑 Sigmoid 激活函数。

\vspace{1em}

\textbf{6.8 ($\star$)} Show that minimization of the error function (6.25) with respect to the variance $\sigma^2$ gives the result (6.27).

\textbf{6.8 ($\star$)} 证明关于方差 $\sigma^2$ 最小化误差函数 (6.25) 会得到结果 (6.27)。

\vspace{1em}

\textbf{6.9 ($\star$)} Show that maximizing the likelihood function under the conditional distribution (6.28) for a multioutput neural network is equivalent to minimizing the sum-of-squares error function (6.29). Also, show that the noise variance that minimizes this error function is given by (6.30).

\textbf{6.9 ($\star$)} 证明在多输出神经网络中，最大化条件分布 (6.28) 下的似然函数等价于最小化平方和误差函数 (6.29)。此外，证明使该误差函数最小化的噪声方差由 (6.30) 给出。

\vspace{1em}

\textbf{6.10 ($\star\star$)} Consider a regression problem involving multiple target variables in which it is assumed that the distribution of the targets, conditioned on the input vector $\mathbf{x}$, is a Gaussian of the form
\begin{equation}
p(\mathbf{t}|\mathbf{x}, \mathbf{w}) = \mathcal{N}(\mathbf{t}|\mathbf{y}(\mathbf{x}, \mathbf{w}), \boldsymbol{\Sigma})
\tag{6.66}
\end{equation}
where $\mathbf{y}(\mathbf{x}, \mathbf{w})$ is the output of a neural network with input vector $\mathbf{x}$ and weight vector $\mathbf{w}$, and $\boldsymbol{\Sigma}$ is the covariance of the assumed Gaussian noise on the targets. Given a set of independent observations of $\mathbf{x}$ and $\mathbf{t}$, write down the error function that must be minimized to find the maximum likelihood solution for $\mathbf{w}$, if we assume that $\boldsymbol{\Sigma}$ is fixed and known. Now assume that $\boldsymbol{\Sigma}$ is also to be determined from the data, and write down an expression for the maximum likelihood solution for $\boldsymbol{\Sigma}$. Note that the optimizations of $\mathbf{w}$ and $\boldsymbol{\Sigma}$ are now coupled, in contrast to the case of independent target variables discussed in Section 6.4.1.

\textbf{6.10 ($\star\star$)} 考虑一个涉及多个目标变量的回归问题，其中假设在给定输入向量 $\mathbf{x}$ 的条件下，目标的分布是高斯分布，形式为
\begin{equation}
p(\mathbf{t}|\mathbf{x}, \mathbf{w}) = \mathcal{N}(\mathbf{t}|\mathbf{y}(\mathbf{x}, \mathbf{w}), \boldsymbol{\Sigma})
\tag{6.66}
\end{equation}
其中 $\mathbf{y}(\mathbf{x}, \mathbf{w})$ 是一个神经网络的输出，其输入向量为 $\mathbf{x}$，权重向量为 $\mathbf{w}$，而 $\boldsymbol{\Sigma}$ 是假设的目标上高斯噪声的协方差。给定一组独立的 $\mathbf{x}$ 和 $\mathbf{t}$ 的观测值，如果我们假设 $\boldsymbol{\Sigma}$ 是固定且已知的，请写出必须最小化以找到 $\mathbf{w}$ 的最大似然解的误差函数。现在假设 $\boldsymbol{\Sigma}$ 也需要从数据中确定，请写出 $\boldsymbol{\Sigma}$ 的最大似然解的表达式。请注意，$\mathbf{w}$ 和 $\boldsymbol{\Sigma}$ 的优化现在是耦合的，这与 6.4.1 节中讨论的目标变量相互独立的情况不同。

\vspace{1em}

\textbf{6.11 ($\star\star$)} Consider a binary classification problem in which the target values are $t \in \{0, 1\}$, with a network output $y(\mathbf{x}, \mathbf{w})$ that represents $p(t=1|\mathbf{x})$, and suppose that there is a probability $\epsilon$ that the class label on a training data point has been incorrectly set. Assuming i.i.d. data, write down the error function corresponding to the negative log likelihood. Verify that the error function (6.33) is obtained when $\epsilon = 0$. Note that this error function makes the model robust to incorrectly labelled data, in contrast to the usual cross-entropy error function.

\textbf{6.11 ($\star\star$)} 考虑一个二分类问题，其中目标值为 $t \in \{0, 1\}$，网络输出 $y(\mathbf{x}, \mathbf{w})$ 代表 $p(t=1|\mathbf{x})$，并假设训练数据点的类别标签被错误设置的概率为 $\epsilon$。假设数据是独立同分布的，请写出对应于负对数似然的误差函数。验证当 $\epsilon = 0$ 时可以得到误差函数 (6.33)。请注意，与通常的交叉熵误差函数相比，这个误差函数使模型对错误标记的数据具有鲁棒性。

\vspace{1em}

\textbf{6.12 ($\star\star$)} The error function (6.33) for binary classification problems was derived for a network having a logistic-sigmoid output activation function, so that $0 \le y(\mathbf{x}, \mathbf{w}) \le 1$, and data having target values $t \in \{0, 1\}$. Derive the corresponding error function if we consider a network having an output $-1 \le y(\mathbf{x}, \mathbf{w}) \le 1$ and target values $t=1$ for class $\mathcal{C}_1$ and $t=-1$ for class $\mathcal{C}_2$. What would be the appropriate choice of output-unit activation function?

\textbf{6.12 ($\star\star$)} 用于二分类问题的误差函数 (6.33) 是为具有逻辑 Sigmoid 输出激活函数的网络推导出来的，因此 $0 \le y(\mathbf{x}, \mathbf{w}) \le 1$，并且数据的目标值为 $t \in \{0, 1\}$。如果我们考虑一个网络的输出为 $-1 \le y(\mathbf{x}, \mathbf{w}) \le 1$，并且对于类别 $\mathcal{C}_1$ 目标值为 $t=1$，对于类别 $\mathcal{C}_2$ 目标值为 $t=-1$，请推导出相应的误差函数。应该选择什么样的输出单元激活函数？

\vspace{1em}

\textbf{6.13 ($\star$)} Show that maximizing the likelihood for a multi-class neural network model in which the network outputs have the interpretation $y_k(\mathbf{x}, \mathbf{w}) = p(t_k=1|\mathbf{x})$ is equivalent to minimizing the cross-entropy error function (6.36).

\textbf{6.13 ($\star$)} 证明对于一个多类神经网络模型，如果其网络输出的解释为 $y_k(\mathbf{x}, \mathbf{w}) = p(t_k=1|\mathbf{x})$，那么最大化其似然等价于最小化交叉熵误差函数 (6.36)。

\vspace{1em}

\textbf{6.14 ($\star$)} Show that the derivative of the error function (6.33) with respect to the pre-activation $a_k$ for an output unit having a logistic-sigmoid activation function $y_k = \sigma(a_k)$, where $\sigma(a)$ is given by (6.13), satisfies (6.31).

\textbf{6.14 ($\star$)} 证明对于具有逻辑 Sigmoid 激活函数 $y_k = \sigma(a_k)$ 的输出单元，其误差函数 (6.33) 关于预激活值 $a_k$ 的导数满足 (6.31)，其中 $\sigma(a)$ 由 (6.13) 给出。

\vspace{1em}

\textbf{6.15 ($\star$)} Show that the derivative of the error function (6.36) with respect to the pre-activation $a_k$ for output units having a softmax activation function (6.37) satisfies (6.31).

\textbf{6.15 ($\star$)} 证明对于具有 Softmax 激活函数 (6.37) 的输出单元，其误差函数 (6.36) 关于预激活值 $a_k$ 的导数满足 (6.31)。

\vspace{1em}

\textbf{6.16 ($\star\star$)} Write down a pair of equations that express the Cartesian coordinates $(x_1, x_2)$ for the robot arm shown in Figure 6.16 in terms of the joint angles $\theta_1$ and $\theta_2$ and the lengths $L_1$ and $L_2$ of the links. Assume the origin of the coordinate system is given by the attachment point of the lower arm. These equations define the forward kinematics of the robot arm.

\textbf{6.16 ($\star\star$)} 写下一对方程，用关节角度 $\theta_1$ 和 $\theta_2$ 以及连杆长度 $L_1$ 和 $L_2$ 来表示图 6.16 所示机器人手臂的笛卡尔坐标 $(x_1, x_2)$。假设坐标系的原点由下臂的连接点给出。这些方程定义了机器人手臂的正向运动学。

\vspace{1em}

\textbf{6.17 ($\star\star$)} Show that the variable $\gamma_{nk}$ defined by (6.44) can be viewed as the posterior probabilities $p(k|\mathbf{t})$ for the components of the mixture distribution (6.38) in which the mixing coefficients $\pi_k(\mathbf{x})$ are viewed as $\mathbf{x}$-dependent prior probabilities $p(k)$.

\textbf{6.17 ($\star\star$)} 证明由 (6.44) 定义的变量 $\gamma_{nk}$ 可以看作是混合分布 (6.38) 中各分量的后验概率 $p(k|\mathbf{t})$，其中混合系数 $\pi_k(\mathbf{x})$ 被视为依赖于 $\mathbf{x}$ 的先验概率 $p(k)$。

\vspace{1em}

\textbf{6.18 ($\star\star$)} Derive the result (6.45) for the derivative of the error function with respect to the network output pre-activations controlling the mixing coefficients in the mixture density network.

\textbf{6.18 ($\star\star$)} 推导出误差函数关于控制混合密度网络中混合系数的网络输出预激活值的导数的结果 (6.45)。

\vspace{1em}

\textbf{6.19 ($\star\star$)} Derive the result (6.46) for the derivative of the error function with respect to the network output pre-activations controlling the component means in the mixture density network.

\textbf{6.19 ($\star\star$)} 推导出误差函数关于控制混合密度网络中分量均值的网络输出预激活值的导数的结果 (6.46)。

\vspace{1em}

\textbf{6.20 ($\star\star$)} Derive the result (6.47) for the derivative of the error function with respect to the network output pre-activations controlling the component variances in the mixture density network.

\textbf{6.20 ($\star\star$)} 推导出误差函数关于控制混合密度网络中分量方差的网络输出预激活值的导数的结果 (6.47)。

\vspace{1em}

\textbf{6.21 ($\star\star\star$)} Verify the results (6.48) and (6.50) for the conditional mean and variance of the mixture density network model.

\textbf{6.21 ($\star\star\star$)} 验证混合密度网络模型的条件均值和方差的结果 (6.48) 和 (6.50)。

\end{document}
